\section{Introduction}
PCPOP is a Julia library that builds and solves SDP relaxations for polynomial optimization with partially-commutative variables, where some variables commute and some not, and hence generalizing the non-commutative and fully-commutative polynomial optimization. It supports both the eigen-value and tracial versions.
%  The algebra associated with Partially-commutative polynomial optimization is the Partially-commutative polynomial algebra[see \cite{}] and it can be seen as special case of Graph product polynomial algebra [see \ref{sec:}]. At the core, 
% \begin{itemize}
%     \item PCPOP implements various data structures to efficiently build monomials belonging to Graph product algebra and do arithmetic with them.  
%     \item PCPOP implements various algorithms, the theory of which is mentioned in \cite{}, to find normal forms of polynomials under given polynomial identities. 
% \end{itemize}


Polynomial optimization has become a central tool in many diverse fields, but it is a NP-hard problem. Lassere first developed a numerical algorithm to solve the fully-commutative polynomial optimization problems. This was later extended to Non-Commutative variables yielding the NPA hierarchy \cite{}. Given the natural non-commutativity of operators in quantum physics, the NPA hierarchy finds numerous applications in quantum foundations, device-independent quantum cryptography, many-body physics, and others. It has also been extended to trace polynomials, with many prominent applications, such as in quantum networks and quantum cryptography in prepare and measure scenarios. 

Given $p(\underline{x}),q_i(\underline{x})$, polynomials in \underline{x}=($x_1,x_2...x_n$) variables, the fully-commutative polynomial optimization reads as,
\begin{equation}
\begin{split}
       p^* \: = \:\underset{\underline{x}}{inf} \: p(\underline{x}), \\
    s.t \: q_i(\underline{x}) \geq 0. 
\end{split}
\end{equation}
In non-commutative variables \textbf{X}=($\textbf{x}_1,\textbf{x}_2...\textbf{x}_n$) and polynomials $p(X),q_i(X), r_i(X)$ in free algebra, the eigen-value version of non-commutative polynomial optimization reads as,
\begin{equation}
  \label{ncpop}
\begin{split}
       p^* =\underset{\underline{A} \in \mathcal{B}(\mathcal{H})^n, \psi \in \mathcal{H}, \mathcal{H}}{inf} \: \braket{\psi,p(\underline{A}) \psi}, \\
    s.t \: q_i(\underline{A}) \geq 0, \\
    s.t \: \braket{\psi ,r_i(\underline{A}) \psi} \geq 0, \\
    s.t \: \braket{\psi , \psi} = 1. \\ 
\end{split}
\end{equation}
In non-commutative variables \textbf{X}=($\textbf{x}_1,\textbf{x}_2...\textbf{x}_n$) and polynomials $p(X),q_i(X), r_j(X)$, in free algebra, the tracial version of non-commutative polynomial optimization reads as,
\begin{equation}
\begin{split}
       p^* =& inf \: tr(p(X)), \\
    s.t \: & q_i(X) \geq 0. \\
     & tr(r_j(X)) \geq 0
\end{split}
\end{equation}
Prtially commutative polynomial optimization is a polynomial optimization in the partially-commuative algebra[see \cite{} for details]. Informally speaking, given set of variables $X$ and a graph $\graph=(X,E)$, whose vertices are letters that are connected if they commute, we denote by $X^\graph$ the partially-commutative variables and the partially commutative polynomial optimization reads as,
\begin{equation}
  \label{pcpop}
\begin{split}
       p^* =\underset{\underline{A} \in \mathcal{B}(\mathcal{H})^n, \psi \in \mathcal{H}, \mathcal{H}}{inf} \: \braket{\psi,\canonical{p}(\underline{A}) \psi}, \\
    s.t \: \canonical{q_i}(\underline{A}) \geq 0, \\
    s.t \: \braket{\psi ,\canonical{r_i}(\underline{A}) \psi} \geq 0, \\
    s.t \: \braket{\psi , \psi} = 1. \\ 
\end{split}
\end{equation}
where $\canonical{p}$ means the polynomial with partially-commutative variables. The tracial version can be defined similarly as,
\begin{equation}
  \label{tpop}
\begin{split}
       p^* =& inf \: tr(\canonical{p}(X)), \\
    s.t \: & \canonical{q_i}(X) \geq 0. \\
     & tr(\canonical{r_j}(X)) \geq 0
\end{split}
\end{equation}

The algorithms for solving polynomial optimization problems, whether commutative or non-commutative, and their tracial or eigenvalue versions, share the same core idea of building a sequence of SDP relaxations of increasing size, whose optimal values converge to the global optimum under certain assumptions.  The SDP relaxations are often indexed by a level parameter, with higher levels yielding better approximations of the global optimum but at the cost of increased computational complexity.

\subsection{Motivation for PCPOP}
Given how intractable it can be to solve the SDPs for higher levels or bigger problem sizes, i.e., with bigger PSD matrices and a large number of PSD and scalar constraints, several solutions to reduce the problem complexity by harnessing the structure of the input problem have been proposed. For e.g in, they use the sparsity of the input problem to reduce the SDP size at each levels of hierarchy, in \cite{} , they propose how to use symmetry in the input problem to reduce the problem size and in \cite{} they propose efficient solution for special type of correlations called synchronous correlations.

The main motivation for PCPOP comes from polynomial optimization problems with equality constraints that appear ubiquitously in physics, especially in the field of quantum information. We found mathematical tools that abstracts these ubiquitous scenarios and PCPOP implements the associated mathematical structures and algorithms to efficiently build and solve the SDP relaxations. Below we discuss three classes of problems that motivated us to build PCPOP\@

For the first class, consider a quantum experiment, where some operators act jointly on quantum systems(say an unitary describing interaction of two quantum systems) and some operators act locally (say measurement on a quantum system). So operators that share a quantum system don't commute with each other but operators that don't share a quantum system , that is, they act on space-like separated  quantum systems commute with each other. The operators could also satisfy other binomial constraints such as being projectors, unitary, dichotomic, pauli constraints or some pairs of operators being orthogonal.  These operators can be abstracted as partially-commutative variables. And the polynomial optimization consisting of such operators can be best solved as a partially-commutative polynomial optimization problem[see \ref{} for examples and how they are implemented]. 


For the second class, consider a polynomial optimization problem which consists of two types of equality constraints. First being, special commutation rules such that the variables can be partitioned into different subsets where all the variables in a subset share the same commutation relations. The other type of constraints consist of polynomials that act on variables belonging to the same subset. We discuss methods in \cite{} to leverage such type of equality constraints as efficiently as possible. The abstract mathematical structure used for this is the Graph product algebra which extends the partially-commutative algebra. See \ref{} for examples and how they are implemented.

The third class of problems consists of trace polynomial optimization problems as in \ref{} but with partially commutative variables and the equality constraints requiring the variables to be projectors, unitaries, or dichotomic. This class of problems appear naturally when considering prepare and measure scenrios. We discuss methods in \cite{} to build efficient SDP relaxations for such class of problems. See \ref{} for examples and how they are implemented.

\subsection{Scope of PCPOP}
Here we discuss all the classes of polynomial optimization problems that can be solved using PCPOP.
\begin{itemize}
  \item Non-commutative polynomial optimization problems, as mentioned in \ref{ncpop} and illustrated in section \ref{sec:}. PCPOP uses AbstractAlgebra.jl package to reduce the polynomials using non-commutative Gr\"obner basis. 
  \item Partially-commutative polynomial optimization problems, as mentioned in \ref{pcpop} and illustrated in section \ref{sec:}. Here the special binomial constraints such as projectors, unitaries, dichotomic, orthogonality  etc. are taken care implicitly while building the monomials as clique normal forms.
  \item Graph product polynomial optimization problems.
  \item Tracial polynomial optimization problems with partially commuting variables, as mentioned in \ref{tpop} and illustrated in \ref{sec:}. Here the special binomial constraints such as projectors, unitaries, dichotomic, orthogonality  etc. are taken care implicitly while building the monomials as cyclic graph normal forms.
\end{itemize}
\subsection{How PCPOP compares to other related libraries}
Ncpol2sdpa is a python-based library for solving non-commutative polynomial optimization problems. It is popular in the quantum information community as it was the first library to implement the NPA hierarchy with easy interface to fully express a non-commutative polynomial optimization problem. Ncpol2sdpa has some shortcomings as mentioned below, that PCPOP tries to overcome.
\begin{itemize}
    \item It is python-based, and hence slower than Julia-based libraries.
    \item It implicitly assumes normalization and it is hard to tweak that.
    \item Trace polynomial optimization[\ref{}] is not implemented.
    \item Support for leveraging symmetries in the optimization is very limited.
    \item General equality constraints are implemented using localizing matrices which is not efficient and binomial constraints are handled using replacement rules which isn't effective [see \ref{} for an example].
\end{itemize}
QuantumNPA is another library for implementing NPA hierarchy. It is Julia-based and has easier interface to express common quantum information scenarios. It also has some short-comings that PCPOP tries to overcome.
\begin{itemize}
  \item Support for leveraging symmetries in the optimization problem is very limited.
  \item Trace polynomial optimization[\ref{}] is not implemented.
  \item Hard to impose general equality constraints.
  \item It can handle partially-commutative varibles that satisfy common binomial constraints such as being projector, unitary or dichotomic but it can be hard to express those and not efficient for larger problems.
\end{itemize}
Another popular Julia-based library is SumOfSquares which is excellent for fully-commutative polynomial optimization problems and has support for leveraging sparsity and symmetry in the input problem. It has some support for non-commutative polynomial optimization problems with shortcomings that PCPOP tries to overcome.
\begin{itemize}
  \item Hard to impose general equality constraints.
  \item Lack of support for quantum information specific interface or simplifications.
  \item Trace polynomial optimization[\ref{}] is not implemented.
  \item It lacks support for leveraging symmetries in the optimization problem for the non-commutative case.
\end{itemize}
% A concrete example will be a Bell scenario with two parties, Alice and Bob, where measurements of Alice or Bob don't commute with each other but Alice's measurements commute with Bob's measurements and one is interested in the highest quantum violation of CHSH expression.
